\subsection{Related work for Part II}\label{sec:lit-part2}

Part II asks which activities a firm should keep in house, outsource, automate, augment or eliminate, and how those choices add up to the firm-level parameters of Part I. We review the antecedents of each component of the model in turn.

\paragraph{Modes and boundaries.} The unit of analysis is the task. \citet{zeira1998workers} modeled technologies that lower labor requirements but raise capital requirements and showed that they are adopted only where productivity is high. \citet{autor2003skill} found that computers substitute for workers in routine cognitive and manual tasks and complement them in nonroutine problem solving, and \citet{acemoglu2011skills} built the task-based framework in which skills are assigned to tasks and machines can take over some of them. \citet{acemoglu2018race} added new tasks in which labor has a comparative advantage, \citet{acemoglu2019automation} separated displacement from reinstatement effects on labor demand, and \citet{acemoglu2022tasks} linked 50 to 70 percent of the change in the U.S. wage structure since 1980 to task displacement, alongside modest productivity gains. Their automation margin, a task-by-task comparison of the costs of capital and labor, is our crossover $c_A=c_H$. The option of buying the task comes from the theory of the firm. \citet{coase1937nature} explained the boundary of the firm by the cost of using the market, \citet{williamson1985economic} made transaction costs, bounded rationality and opportunism the basis for choosing between market and hierarchy, and \citet{grossman1986costs} defined ownership as the residual rights of control that matter when contracts are incomplete. Our coordination share $\kappa_O$, risk premium $\zeta_O$ and separability condition are reduced forms of these arguments. \citet{grossman2008trading} modeled offshoring as trade in tasks and showed that cheaper task trade raises the productivity of the factor whose tasks move, so the outsourcing price $p_O$ moves with the cost of trading tasks as well as with vendors' adoption of AI.

\paragraph{Complementarity and integration.} \citet{kremer1993oring} wrote a production function in which a mistake in any one of several tasks can destroy the output, so that quality multiplies across tasks; our process success probability $\prod_j r_j$ is his. \citet{milgrom1990economics} modeled the joint adoption of flexible technology and new organizational forms, with comparative statics that require neither differentiability nor convexity, and \citet{topkis1998supermodularity} and \citet{milgrom1994monotone} supply the lattice methods. Supermodularity induced by handoff costs is what makes our optimal mode vector all or nothing on a process. Organizational investment strongly shapes the return to computers \citep{brynjolfsson2000beyond}, and \citet{brynjolfsson2021productivity} showed that intangible complements to a general purpose technology produce a productivity J-curve, with measured productivity understated early and overstated later. Data-driven decision making nearly tripled in U.S. manufacturing between 2005 and 2010 \citep{brynjolfsson2016rapid}; predictive analytics raised plant performance only where at least one complement was present, namely accumulated IT capital, educated workers or high flow-efficiency production \citep{brynjolfsson2021power}; and early AI adoption, by fewer than 6 percent of U.S. firms in 2018, often occurred together with cloud computing \citep{mcelheran2024ai}. \citet{jones2020nonrivalry} and \citet{farboodi2021model} model data as a nonrival input and as an intangible asset that depreciates. Our readiness index $\iota$ and threshold $\iota^*$ formalize this complement for AI.

\paragraph{Organization.} \citet{garicano2000hierarchies} modeled hierarchies as a way of organizing knowledge: production workers solve common problems and pass exceptions to specialized solvers, so the costs of acquiring and of communicating knowledge determine spans and layers. Cheaper knowledge acquisition and cheaper communication have opposite effects on organization and inequality \citep{garicano2006organization,garicano2015knowledge}, and \citet{bloom2014distinct} confirmed the distinction in firm data: information technologies widen autonomy and spans of control, while communication technologies reduce autonomy. \citet{ide2025artificial} placed AI agents in a knowledge hierarchy and found that autonomous AI mainly benefits the most knowledgeable people and non-autonomous copilots the least knowledgeable. \citet{agrawal2019artificial} treat AI as prediction, an input to decisions that it may either automate or enhance, and \citet{agrawal2022power} argue that the larger gains require redesigning systems of decisions. On team size, \citet{becker1992division} argued that specialization is limited mainly by the cost of coordinating specialists and by the stock of general knowledge, not by the extent of the market, and \citet{brooks1975mythical} observed that adding people to a late software project makes it later. \citet{rosen1981economics} showed how technologies that let one person serve a whole market concentrate output and earnings among the most talented, the mechanism behind a solo contributor who directs agents. In a field experiment with 776 professionals at Procter \& Gamble, individuals working with AI matched two-person teams working without it \citep{dellacqua2026cybernetic}.

\paragraph{Capability and cost dynamics.} We measure capability as a task horizon. \citet{kwa2025measuring} define the 50 percent time horizon, the human completion time of tasks that an AI agent completes half the time, and estimate that it doubled about every seven months (207 days) between 2019 and early 2025, faster for models released since 2024; 80 percent horizons are four to six times shorter but grow at a similar rate. Their tasks are software tasks, and models do worse on ``messier'' tasks of equal length, so applying the horizon to other work is an assumption. Usage data give different levels: for a model with a horizon of about 2 hours on METR's tasks, \citet{appel2026anthropic} estimate about 3.5 hours on API traffic and about 19 hours on Claude.ai, partly because users break tasks down and select them. On prices, the cost of reaching a fixed benchmark score fell by 9 to 900 times a year, depending on the score \citep{epoch2025inference}, and by about 5 to 10 times a year at fixed performance in the estimates of \citet{gundlach2025price}, who also find that the cost of running the record-performance model on each benchmark rose by 3 to 18 times a year. Energy is a small share of cost: servers account for about 60 percent of the annual cost of a one-gigawatt AI data center and energy for about \$0.6 billion of \$8.5 billion \citep{michael2026servers}, and the median Gemini text prompt used 0.24 watt-hours after a 33-fold fall in energy per prompt over one year \citep{elsworth2025measuring}. Falling costs give adoption an option value. Irreversible investment under uncertainty should wait until benefits exceed costs by a margin \citep{mcdonald1986value,dixit1994investment}, and when better technologies arrive at random the optimal pace of adoption is slower than the net present value rule implies \citep{farzin1998optimal}. In the aggregate, \citet{korinek2024scenarios} model the transition to AGI as the automation of all tasks, with wages set by a race between automation and capital accumulation, while \citet{acemoglu2025simple} expects modest productivity effects from current AI over the next decade.

\paragraph{Taxes, regulation and preferences.} \citet{acemoglu2020does} argue that the U.S. tax system taxes labor more heavily than capital and so pushes automation beyond its socially optimal level, because marginal automated tasks add little productivity but displace workers. \citet{guerreiro2022should} find that a robot tax is optimal while routine workers who cannot switch occupations remain in the labor force, and zero afterward, and \citet{costinot2023robots} derive sufficient-statistic formulas for taxes on technologies that create winners and losers. These papers treat tax wedges as policy choices; we treat them as parameters of the firm's problem. Regulation enters through human oversight. Article 14 of the EU AI Act requires high-risk systems to be designed so that natural persons can oversee them, decide not to use them, override their output and stop them, and requires confirmation by at least two people before acting on a remote biometric identification \citep{eu2024regulation}, and Annex III lists the high-risk areas, among them employment, education, credit scoring and the administration of justice. Article 22 of the GDPR restricts decisions based solely on automated processing that significantly affect a person and, where contract or consent allows them, guarantees a right to human intervention \citep{eu2016regulation}. Neither text sets a minimum share of work to be reviewed, so our floor $\psi_{\min}$ is a modeling device that converts these obligations into cost. Preferences enter on the demand side. People avoid algorithms after seeing them err, even when the algorithms outperform humans \citep{dietvorst2015algorithm}; patients resist medical AI because they believe it neglects their individual characteristics \citep{longoni2019resistance}; and in a field experiment, disclosing that a sales call came from a chatbot reduced purchases by more than 79 percent, although undisclosed chatbots sold as well as skilled staff \citep{luo2019machines}; these results motivate the demand penalty $\delta_H$.

\paragraph{Measurement of tasks.} The O*NET content model describes occupations with cross-job descriptors arranged in a hierarchical taxonomy \citep{peterson2001understanding}; the current database links occupational task statements to about 2,100 Detailed Work Activities, each nested in an Intermediate Work Activity \citep{onet2026database}; and occupations are coded to the 867 detailed occupations of the 2018 Standard Occupational Classification \citep{bls2017soc}. For processes, the APQC Process Classification Framework arranges more than 1,000 processes in 13 enterprise-level categories \citep{apqc2024process}. Measures of AI exposure use these units: the patent-text measure of \citet{webb2020impact}, which finds that AI, unlike software and robots, targets high-skilled tasks; the suitability-for-machine-learning rubric of \citet{brynjolfsson2018what}; the occupational exposure index of \citet{felten2021occupational}; and the exposure ratings for large language models of \citet{eloundou2024gpts}. Usage data map conversations to the same tasks: about 36 percent of occupations used AI for at least a quarter of their tasks, with 57 percent of use augmenting and 43 percent automating \citep{handa2025which}; directive delegation rose from 27 to 39 percent in eight months, and business use through the API is mostly automation \citep{appel2025anthropic}; and information work is the most common and most successful kind of AI-assisted activity \citep{tomlinson2025working}. Field experiments bound the augmentation gain $\beta_U$. Professionals completed writing tasks in about 40 percent less time with ChatGPT \citep{noy2023experimental}; customer-support agents gained most when least experienced \citep{brynjolfsson2025generative}; developers finished a coding task 55.8 percent faster with Copilot \citep{peng2023impact}, and pooled experiments at three firms found 26 percent more completed tasks \citep{cui2024effects}; consultants completed 12.2 percent more tasks, 25.1 percent faster, inside the frontier of AI capability but were less often correct on a task outside it \citep{dellacqua2026navigating}; and experienced open-source developers took 19 percent longer with early-2025 tools \citep{becker2025measuring}. Employment effects are so far small overall \citep{acemoglu2022artificial,humlum2025still} but concentrated among young workers in exposed occupations where AI automates rather than augments \citep{brynjolfsson2025canaries}. \citet{howard1966information} supplies the value of information we use to price finer cost data.

\paragraph{Transitions, risk and decision capacity.} Moving an activity between modes is itself a risky project. Across 1,471 IT projects, the average cost overrun was 27 percent, but one project in six overran its budget by 200 percent on average and its schedule by almost 70 percent \citep{flyvbjerg2011why}; across 5,392 projects, cost overruns follow a power-law distribution with a tail exponent of about 2.3, so an assumption of normal-like overruns understates extreme outcomes \citep{flyvbjerg2022empirical}. Adoption is followed by a dip. After a manufacturer switched on an enterprise system, lead times and on-time delivery worsened before improving beyond their earlier levels \citep{mcafee2002impact}; ERP adopters show a slowdown in performance and productivity shortly after implementation but a higher market valuation \citep{hitt2002investment}; the measured contribution of computers is up to five times larger over five to seven years than over one, because complementary investment takes time \citep{brynjolfsson2003computing}; and whether an ERP project counts as a success depends on when it is measured \citep{markus2000learning}. These results motivate our parallel-run costs, ramp and rollback probability. Pilots and stage gates are options. Funding rules for sequential projects must allow for later reviews \citep{roberts1981funding}, the optimal search among uncertain alternatives is a reservation-price rule \citep{weitzman1979optimal}, a simple net present value rule can lead to gross over-investment in projects that take time to build \citep{majd1987time}, and deferral, expansion, contraction and abandonment can all be valued as options \citep{trigeorgis1996real,dixit1994investment}. Learning across waves is real but perishable: knowledge acquired through production depreciates quickly and transfers imperfectly \citep{argote1990persistence,darr1995acquisition}, and in software services, team familiarity and role experience predict performance \citep{huckman2009team}. The capacity to decide is limited. \citet{simon1947administrative} built administrative theory on bounded rationality, \citet{simon1971designing} observed that a wealth of information creates a poverty of attention, and \citet{ocasio1997towards} explains firm behavior by how the firm channels its decision makers' attention. Formal models treat managers as processors of limited capacity, so that organizations trade the number of processors against delay \citep{radner1993organization,vanzandt1999realtime}, managers choose what to know \citep{geanakoplos1991theory}, and scarce attention forces a firm to coordinate only a few tasks \citep{dessein2016rational}. Hierarchies reject more good projects and polyarchies accept more bad ones \citep{sah1986architecture}, which links decision rights to our success probability, and how chief executives spend their time is related to firm performance \citep{bandiera2020ceo}. Finally, labor adjustment is costly and poorly described by symmetric quadratic costs \citep{hamermesh1996adjustment}, severance requirements reduce employment \citep{lazear1990job}, and high-tenure workers displaced from distressed firms lose about 25 percent of earnings a year over the long term \citep{jacobson1993earnings}, which bears on both the severance term and the altruism wedge.

\paragraph{Leadership change, reorganization and organizational memory.} Transformations outlast many of the people who start them. Annual CEO turnover at large U.S. firms was 15.8 percent over 1992 to 2007, implying an average tenure of under seven years, and 16.8 percent after 2000 \citep{kaplan2012how}, and boards dismiss CEOs after industry and market shocks outside their control \citep{jenter2015ceo}. Leaders shape what follows: manager fixed effects explain a substantial share of differences in investment, financial and organizational practices \citep{bertrand2003managing}, and a successor makes less strategic change when the predecessor stays on as board chair \citep{quigley2012when}. These results inform the sponsor hazard $h_L$ and the retention probability $\kappa_{\mathrm{keep}}$. Reorganizations are disruptive in their own right: each change raised the immediate risk of failure among Finnish newspapers \citep{amburgey1993resetting}, and among large U.S. industrial firms, reconfiguring units without changing design principles tended to lower performance, except in dynamic environments \citep{girod2017reconfiguration}. \citet{walsh1991organizational} found earlier treatments of organizational memory fragmented and developed a more coherent account of it, and losing the people who hold that memory is costly. Turnover rates are negatively related to organizational performance \citep{hausknecht2011collective,park2013turnover}; an inventor's premature death persistently lowers coinventors' earnings and patents \citep{jaravel2018team}; and in retail stores, turnover hurts margins and service only where managers enforce little process conformance \citep{ton2008managing}, the insurance role that our documentation parameter $\kappa_{\mathrm{doc}}$ plays. A new sponsor also inherits sunk commitments. Decision makers escalate commitment to failing courses of action for which they are responsible \citep{staw1976knee,staw1981escalation,brockner1992escalation}, and 30 to 40 percent of information-system projects show some escalation \citep{keil2000why}. The counterweight is the abandonment option, which investors price: firm value rises with liquidation value, more so where assets are general-purpose \citep{berger1996investor}, which supports our preference for vision-agnostic investments with a high $\vartheta_{\mathrm{reuse}}$.

\paragraph{Ownership and horizons.} The owner sets a transformation's discount rate, horizon and constraints. Average private equity fund returns net of fees roughly matched the S\&P 500, with wide dispersion and persistence across partnerships \citep{kaplan2005private}; with better data, U.S. buyout funds beat the S\&P 500 by 20 to 27 percent over a fund's life, more than 3 percent a year, while venture funds did better in the 1990s and worse in the 2000s \citep{harris2014private}. The median buyout is held about six years, and investors add value through financial, governance and operational engineering \citep{kaplan2009leveraged}. In public-to-private buyouts of 1990 to 2006, operating gains were comparable to or slightly above those of benchmark firms and returns tracked rising industry multiples \citep{guo2011do}; after going private, plant productivity gains relative to peers are small but plant portfolios are restructured faster \citep{bharath2014do}; and restaurants became cleaner and safer after buyouts, especially when the partners had industry experience \citep{bernstein2016operational}. Public firms face short-term pressure. Managers forgo good investments to raise current earnings even when the market sees through it \citep{stein1989efficient}; in a survey, 78 percent of executives would give up economic value to smooth earnings and 55 percent would forgo a positive-NPV project rather than miss the quarter's consensus \citep{graham2005economic}; and public firms invest less, and respond less to investment opportunities, than comparable private firms \citep{asker2015corporate}. Family ownership adds value when the founder is chief executive or chair, not when a descendant is \citep{villalonga2006how}. Leverage binds through covenants, after whose violation investment falls sharply \citep{chava2008how}. For venture capital, estimates of risk and return depend on correcting for the selection of observed valuations \citep{cochrane2005risk,korteweg2010risk}. These results bound the regime presets: the myopia weight, the hold period and exit channel, and the covenant cap on the dip. Required hurdle rates come only from practitioner targets, so we treat them as assumptions.

\paragraph{Product strategy.} \citet{shapiro1999information} showed how to price and version information and compete under lock-in and network effects. When metering has transaction costs, adding a fixed fee to usage-based pricing is always profitable \citep{sundararajan2004nonlinear}, and bundling many information goods can raise profit per good by making valuations predictable \citep{bakos1999bundling}. Neither subscription nor metering dominates. Two-part tariffs are not always optimal for access services, and with enough spare capacity only flat rates survive \citep{essegaier2002pricing}; pay-per-use beats selling for a monopolist when the meter's psychological cost is low, but not against a rival that sells \citep{balasubramanian2015pricing}; and customers who overpay on flat rates do not churn more, whereas pay-per-use customers who underestimate their usage do \citep{lambrecht2006paying}. Value-based pricing has a decision framework but little academic study \citep{hinterhuber2004towards}; our capture share $\theta_{\mathrm{cap}}$ is its reduced form. Adoption follows \citet{bass1969new}; a meta-analysis of 213 applications related its parameters to the innovation and country \citep{sultan1990meta}, and for 52 consumer durables in 28 countries the mean coefficients of innovation and imitation were 0.027 and 0.419 \citep{vandenbulte2004social}. A monopolist gains less from an innovation than a competitive entrant because it replaces its own profits \citep{arrow1962economic}, and staggered introductions beat simultaneous ones when cannibalization is severe and customers are more impatient than the seller \citep{moorthy1992market}. For AI, a model provider's optimal menu is a committed-spend plan with tokens priced at marginal cost \citep{bergemann2025menu}; scale and scope economies may concentrate the foundation-model market \citep{korinek2025concentrating}; market power depends on whether data can be traded across firms \citep{gans2024market}; and applications multihome across models, and demand for a new model levels off within weeks \citep{fradkin2025demand}. Data can entrench large firms or let data-savvy entrants overtake them \citep{farboodi2019big,farboodi2021model}, depending on how firms learn from data \citep{hagiu2023data}; at Amazon, longer histories improved forecasts but more products did not \citep{bajari2019impact}.

\paragraph{Security, maintenance and ancillary costs.} Because distributed systems are built from machines whose owners have conflicting interests, incentives matter as much as technical design for security \citep{anderson2006economics}. \citet{gordon2002economics} derive optimal security investment: for two broad classes of breach probability functions it never exceeds $1/e$, about 37 percent, of the expected loss, and it does not always rise with vulnerability. The bound is not general: under other breach functions investment can exceed it and jump from zero to its maximum \citep{hausken2006returns}. We therefore treat it as one case for the security investment $s_{\mathrm{sec}}$, not a rule. Agents add an attack surface. Indirect prompt injection, which plants instructions in data that a model later retrieves, let researchers steal data, change behavior and trigger other API calls in deployed applications, against which effective mitigations were lacking \citep{greshake2023not}. Prompt injection heads the practitioner list of risks for LLM applications, which also includes excessive agency and unbounded consumption \citep{owasp2024top}; hence our exposure terms for external inputs and write rights. For governance, the voluntary NIST framework organizes AI risk management into four functions, govern, map, measure and manage \citep{nist2023artificial}, and ISO/IEC 42001 sets requirements for an AI management system \citep{iso2023information}; we treat compliance as a fixed cost. Machine learning systems incur large ongoing maintenance costs through entanglement, hidden feedback loops, data dependencies and a changing external world \citep{sculley2015hidden}, and more complex code costs more to maintain \citep{banker1993software}. Evidence on maintenance's share of cost is weaker than often claimed. A survey of 69 organizations found that maintenance and enhancement consume much of systems and programming resources \citep{lientz1978characteristics}, and a later survey covered 487 organizations \citep{lientz1980software}. Yet maintenance measured on three products, including the Linux kernel and GCC, had more than twice the survey's corrective share \citep{schach2003determining}. The maintenance rate $\mu_m$ and the revalidation cost $C_{\mathrm{reval}}$ therefore remain assumptions.

\paragraph{What Part II adds.} Part II combines these strands in a single firm's decision problem. The task-based literature compares labor with machines and the theory of the firm compares making with buying; we put in-house work, outsourcing, automation, augmentation and elimination on one cost frontier, so that the value of automation is measured against the cheaper of in-house and outsourced work and vendors' own adoption of AI makes timing part of the choice. We derive conditions for all-or-nothing automation of a process from handoff costs, and a data-readiness threshold below which only augmentation pays, which gives the J-curve a task-level source. We relate team size and hierarchy to AI leverage and coordination costs, treat taxes, oversight rules and preferences as parameters of a nested set of activities available to automate, model each transition as a staged, risky project that draws on limited decision capacity and is exposed to leadership, reorganization and ownership shocks, and aggregate activity choices into the parameters of Part I so that a strategy can be scored. The model organizes the estimates reviewed above rather than replacing them: most of its parameters come from this literature, and its outputs inherit their uncertainty, in particular the transfer of horizons measured on software tasks to other work.
